% !TeX root = ../mobile_sched.tex
\chapter{引言}

\section{研究背景与问题定义}

\subsection{移动端本地推理的现实需求}

近年来，大语言模型和多模态模型逐步从云侧走向端侧。手机、平板和便携式边缘设备开始承担更复杂的推理任务。本地推理的价值并不抽象，主要体现在三个方面。第一，数据不出端可以减少隐私泄露风险。第二，弱网和离线场景仍然可以维持基本可用性。第三，交互链路缩短后，系统更容易控制首 token 时延和连续生成时延。对终端产品而言，这三点都比单纯追求峰值吞吐更重要。

与服务器不同，移动端的资源条件十分苛刻。内存容量有限，持续散热能力弱，操作系统调度策略也更偏向综合体验而不是单任务最优。在这样的环境下，CPU 仍然是最稳定、最通用的推理后端。一方面，CPU 不依赖专用驱动和额外框架，部署路径最短；另一方面，CPU 能覆盖更多机型，便于把推理能力落到真实产品中。MNN 之类的轻量推理框架因此具有很强的工程价值\cite{jiang2020mnn}。

但“能跑”不等于“跑得好”。移动端 CPU 往往采用异构多核结构，不同核心在频率、微架构和缓存层次上都存在明显差异。若仍沿用面向对称多核的线程划分方式，线程迁移、负载失衡和尾部阻塞就会频繁出现。对于交互式大模型推理，这些问题会直接表现为首 token 响应变慢、连续输出抖动增大，以及高负载场景下的性能不稳定。

\subsection{异构多核 CPU 上的执行复杂性}

现代手机 SoC 普遍采用 big.LITTLE 或更细粒度的多簇异构设计。大核追求高性能，中核兼顾性能和功耗，小核则更强调后台负载承载能力。对于同一个推理任务，不同核心完成相同计算量所需时间并不一致。若任务划分仍假定各线程处理速度近似相同，结果往往不是“所有线程同时结束”，而是少数线程拖尾、部分大核提前空闲。

除核心性能差异外，线程绑核与线程迁移也会显著影响执行结果。如果不做绑核，线程可能在多个核心之间切换。线程一旦迁移，缓存热数据会失效，局部性被打断，调度器还会引入额外切换开销。对于持续迭代执行的推理任务，这种开销会反复累积。在大模型推理中，单次额外开销也许不大，但在 Prefill 和 Decode 的多轮执行中会不断放大。

异构多核还带来另一个更隐蔽的问题：最佳线程数并不是固定常数。线程数过少，大核资源无法充分利用；线程数过多，又会引入同步、争用和调度开销。在真实手机上，这条关系常常表现为 U 型曲线。更重要的是，Prefill 和 Decode 的最优线程数并不相同。如果系统强行使用统一线程数和统一绑核策略，至少有一个阶段会偏离最优点。

\subsection{MNN 默认执行方式下的工程瓶颈}

本文研究建立在 MNN 框架之上。MNN 提供了较轻量的部署链路和较完整的 CPU 后端，是移动端本地推理中很有代表性的工业级开源框架\cite{jiang2020mnn}。然而，在面向异构多核 CPU 的大模型推理场景下，MNN 的默认执行路径仍存在进一步优化空间。

首先，默认场景下并不会天然得到适合当前任务的绑核策略。若应用层没有显式指定 affinity，线程可能在运行过程中发生迁移，造成额外调度成本。其次，现有任务划分对核心性能差异的刻画仍然偏粗。固定性能比在某些模型和某些机型上可以工作，但当模型算子形态变化、输入规模变化或系统负载变化时，这种固定参数很容易失准。最后，在高负载环境下，若其他应用抢占部分核心，纯静态划分很难自我纠偏，尾部线程会阻塞整个推理流程。

本课题的直接动机并不是重新提出一套全新的调度理论，而是在真实手机、真实框架和无 root 约束下，把这些工程瓶颈逐项拆开，并给出可复现的改进路径。围绕这一目标，本文把研究问题收束到三个具体方向：分阶段线程数与绑核搜索、基于真实性能比的加权初始划分，以及 Prefill 阶段的无锁双端队列 work stealing。

\section{研究动机与核心问题}

\subsection{绑核与线程数选择对性能的影响}

在本课题的前期实验中，一个非常直接的现象是：绑核后的性能通常优于不绑核。这一现象并不难理解。绑定核心后，线程迁移次数下降，缓存局部性更稳定，运行时的调度扰动也更小。对于需要持续执行大量张量计算的推理任务，这些因素会共同转化为更短的执行时延。

但仅仅“绑核”还不够。继续沿着线程数维度观察，可以看到 Prefill 和 Decode 的性能都不是单调随线程数上升。线程数过小，异构核资源没有被充分利用；线程数继续增大后，划分粒度、同步开销和内存带宽争用开始抬头，性能又会回落。这使得“最佳线程数”成为一个必须通过实测确定的机型相关、阶段相关参数。

进一步地，Prefill 和 Decode 的最优点往往并不重合。Prefill 更像长任务批处理，强调总体吞吐和负载均衡；Decode 更像短步迭代，强调轻量调度和稳定延迟。若继续沿用统一线程数和统一 CPU 集合，系统很难同时兼顾两个阶段。因此，分阶段独立配置线程数和绑核掩码，成为本文系统设计的第一条主线。

\subsection{固定性能比静态划分的局限}

异构多核任务划分的核心问题是“如何估计不同线程的处理能力”。若能力估计不准，初始划分就会失衡。在 MNN 的现有实现中，针对计算密集型任务采用固定性能比做静态分段是一种低成本方案，但它默认了不同机型、不同负载和不同模型下的核心相对速度具有稳定近似值。这个假设在真实手机环境中往往过强。

第一，核心的实际处理能力会受到 DVFS、系统温度和背景任务影响。第二，不同模型的算子组合并不相同，卷积主导、矩阵乘主导和访存主导的负载，对大小核的利用方式也不同。第三，即使同一模型，在 Prefill 和 Decode 两阶段中，算子形态和数据规模也会变化。固定性能比因此很容易出现“平均上看合理，具体到单次运行却不准”的问题。

这意味着，静态划分本身并不是问题，问题在于静态划分之前的能力估计必须更贴近真实设备。本文的第二条主线就是通过实测搜索得到更准确的核心性能比，用真实性能比替代固定经验值，再据此生成加权初始划分边界。

\subsection{高负载场景下的阻塞与长尾问题}

即便性能比估计更准确，纯静态划分仍然存在天然局限：它默认执行过程中不会出现明显扰动。这个假设在实验室空载环境下也许基本成立，但在真实手机上并不稳固。系统服务、界面刷新、后台同步和其他应用都可能竞争 CPU 时间片。当部分核心临时被抢占时，原先看似合理的静态任务边界会迅速变成负担。

Prefill 阶段尤其容易暴露这一问题。该阶段总任务量较大，只要有一两个线程被明显拖慢，整个阶段结束时间就会被尾部线程锁住。换句话说，系统瓶颈不再是平均吞吐，而是最慢线程的完成时间。为此，本文引入基于无锁双端队列的 work stealing 机制：每个线程先执行自己的本地任务；当本地队列为空后，再尝试从其他线程任务尾部窃取剩余工作，以减轻拖尾。

与此相对，Decode 阶段的单步任务更短，调度机制不能太重，否则调度本身就会吃掉收益。因此，本文并不强求两阶段使用同一套策略，而是将 Prefill 设计为“加权初始划分 + work stealing”，将 Decode 设计为更轻量的动态任务获取。这种阶段感知设计，是本文区别于统一策略方案的关键。

\section{本文工作与贡献}

\subsection{分阶段绑核与线程数启发式搜索}

本文首先在真实手机平台上验证了绑核优于不绑核这一基本事实，并进一步观察到 Prefill 和 Decode 的性能随线程数增长呈现明显的 U 型关系。基于这一现象，本文在 MNN 推理流程中引入分阶段的线程数与 affinity 搜索机制，使 Prefill 和 Decode 可以分别使用不同的线程数和不同的 CPU 集合。

\subsection{基于真实性能比的加权初始划分}

针对固定性能比带来的初始任务划分偏差，本文通过实测搜索构建核心性能比，并据此生成加权初始划分边界。该方法不依赖抽象经验参数，而是直接以目标机型的实际执行表现为依据，从而使初始分工更贴近真实异构能力。

\subsection{基于无锁双端队列的 Prefill work stealing 机制}

针对高负载场景下纯静态划分容易产生长尾的问题，本文在 Prefill 阶段设计并实现了基于无锁双端队列的任务窃取机制。各线程先无锁执行本地任务，当本地任务耗尽时，再从其他线程任务尾部窃取剩余工作。该机制保持了静态划分的低开销优势，同时增加了对扰动和失衡的自适应修正能力\cite{blumofe1999workstealing,chase2005deque,le2013weakmemory}。

\section{论文组织结构}

\subsection{章节安排}

第 1 章说明研究背景、工程问题和本文贡献。第 2 章介绍移动端异构多核 CPU、MNN 执行路径、Prefill 与 Decode 的阶段差异，并梳理 work stealing 与移动端推理调度相关工作。第 3 章给出系统设计与关键算法，包括分阶段绑核搜索、性能比标定、加权初始划分和 Prefill work stealing 实现。第 4 章在真实手机平台上完成性能与稳定性评估，并分析高负载场景下的调度行为。第 5 章总结全文并讨论后续工作。

\subsection{写作定位说明}

本文的重点是工程适配与系统验证，而不是构造新的调度理论。全文将尽量围绕“发现了什么瓶颈、为什么这样改、改动如何落到 MNN 中、实验结果说明了什么”这一主线展开，使论文论证与实际实现保持一致。
